% preamble.tex — Asymptotic Theory for Econometricians (White, Revised Edition)
\documentclass[10pt,openright]{book}

%% ── Page Geometry ──────────────────────────────────────────────
% Academic Press format, approximately 6" x 9"
\usepackage[
  paperwidth=6in,
  paperheight=9in,
  top=0.9in,
  bottom=0.9in,
  left=0.75in,
  right=0.75in,
  headheight=14pt,
  headsep=0.25in,
  footskip=0.35in
]{geometry}

%% ── Encoding & Fonts ──────────────────────────────────────────
\usepackage[T1]{fontenc}
\usepackage[utf8]{inputenc}
\usepackage{lmodern}
\usepackage{microtype}

%% ── Math ───────────────────────────────────────────────────────
\usepackage{amsmath}
\usepackage{amssymb}
\usepackage{amsthm}
\usepackage{mathtools}
\usepackage{bm}

% Equation numbering: (chapter.number) — no section prefix
\numberwithin{equation}{chapter}

%% ── Theorem Environments ───────────────────────────────────────
% White uses a single shared counter for ALL theorem-like environments:
% Definition 2.1, Example 2.2, Proposition 2.3, ... all sequential within chapter
\theoremstyle{plain}
\newtheorem{theorem}{Theorem}[chapter]
\newtheorem{lemma}[theorem]{Lemma}
\newtheorem{corollary}[theorem]{Corollary}
\newtheorem{proposition}[theorem]{Proposition}

\theoremstyle{definition}
\newtheorem{definition}[theorem]{Definition}
\newtheorem{example}[theorem]{Example}
\newtheorem{assumption}[theorem]{Assumption}

\theoremstyle{remark}
\newtheorem{remark}[theorem]{Remark}

% QED symbol for proofs
\renewcommand{\qedsymbol}{$\blacksquare$}

%% ── Exercises ──────────────────────────────────────────────────
\usepackage{enumitem}

% White's exercise style: numbered sequentially within chapter, no section prefix
% e.g., Exercise 1, Exercise 2, ... within each chapter
\newcounter{exercise}[chapter]
\renewcommand{\theexercise}{\arabic{exercise}}

\newenvironment{exercises}{%
  \subsection*{Exercises}%
  \begin{enumerate}[label={\textbf{\arabic*.}},ref={\thechapter.\arabic*},leftmargin=*]%
  \setcounter{enumi}{\value{exercise}}%
}{%
  \setcounter{exercise}{\value{enumi}}%
  \end{enumerate}%
}

% Starred exercise marker
\newcommand{\starred}{\textsuperscript{*}}

%% ── Figures ────────────────────────────────────────────────────
\usepackage{graphicx}
\graphicspath{{figures/}{../figures/}}
\usepackage{float}

%% ── Tables ─────────────────────────────────────────────────────
\usepackage{booktabs}
\usepackage{array}
\usepackage{multirow}

%% ── Headers & Footers ──────────────────────────────────────────
\usepackage{fancyhdr}
\pagestyle{fancy}
\fancyhf{}
\fancyhead[LE]{\thepage}
\fancyhead[RE]{\footnotesize\itshape\leftmark}
\fancyhead[RO]{\thepage}
\fancyhead[LO]{\footnotesize\itshape\rightmark}
\renewcommand{\headrulewidth}{0pt}

\renewcommand{\chaptermark}[1]{\markboth{\thechapter.\ #1}{}}
\renewcommand{\sectionmark}[1]{\markright{\thesection\ #1}}

\fancypagestyle{plain}{%
  \fancyhf{}%
  \fancyfoot[C]{\thepage}%
  \renewcommand{\headrulewidth}{0pt}%
}

%% ── TOC Depth ──────────────────────────────────────────────────
% Source TOC shows chapters and sections (no subsections)
\setcounter{tocdepth}{1}

%% ── Cross-references ──────────────────────────────────────────
\usepackage{hyperref}
\hypersetup{
  colorlinks=false,
  pdfborder={0 0 0},
  pdftitle={Asymptotic Theory for Econometricians},
  pdfauthor={Halbert White}
}
\usepackage{cleveref}

%% ── Index ──────────────────────────────────────────────────────
\usepackage{makeidx}
\makeindex

%% ── Custom Commands ────────────────────────────────────────────

% Bold vectors and matrices (White's convention)
\renewcommand{\vec}[1]{\mathbf{#1}}

% Common econometric operators
\DeclareMathOperator{\plim}{plim}
\DeclareMathOperator{\Var}{Var}
\DeclareMathOperator{\var}{var}
\DeclareMathOperator{\Cov}{Cov}
\DeclareMathOperator{\E}{E}
\DeclareMathOperator{\tr}{tr}
\DeclareMathOperator{\rank}{rank}
\DeclareMathOperator{\diag}{diag}
\DeclareMathOperator{\vect}{vec}
\DeclareMathOperator{\sgn}{sgn}
\DeclareMathOperator{\op}{o_p}
\DeclareMathOperator{\Op}{O_p}

% Convergence arrows
\newcommand{\convp}{\xrightarrow{p}}        % convergence in probability
\newcommand{\convd}{\xrightarrow{d}}        % convergence in distribution
\newcommand{\convas}{\xrightarrow{a.s.}}    % almost sure convergence
\newcommand{\convr}[1]{\xrightarrow{r.#1}}  % convergence in r-th mean

% Partial derivatives
\newcommand{\pd}[2]{\frac{\partial #1}{\partial #2}}

% Common notation — bold vectors/matrices
\newcommand{\bfX}{\mathbf{X}}
\newcommand{\bfY}{\mathbf{Y}}
\newcommand{\bfZ}{\mathbf{Z}}
\newcommand{\bfI}{\mathbf{I}}
\newcommand{\bfC}{\mathbf{C}}
\newcommand{\bfW}{\mathbf{W}}
\newcommand{\bfV}{\mathbf{V}}
\newcommand{\bfA}{\mathbf{A}}
\newcommand{\bfB}{\mathbf{B}}
\newcommand{\bfD}{\mathbf{D}}
\newcommand{\bfH}{\mathbf{H}}
\newcommand{\bfM}{\mathbf{M}}
\newcommand{\bfP}{\mathbf{P}}
\newcommand{\bfQ}{\mathbf{Q}}
\newcommand{\bfR}{\mathbf{R}}
\newcommand{\bfS}{\mathbf{S}}
\newcommand{\bfT}{\mathbf{T}}
\newcommand{\bfU}{\mathbf{U}}
\newcommand{\bfbeta}{\boldsymbol{\beta}}
\newcommand{\bfalpha}{\boldsymbol{\alpha}}
\newcommand{\bfgamma}{\boldsymbol{\gamma}}
\newcommand{\bfdelta}{\boldsymbol{\delta}}
\newcommand{\bftheta}{\boldsymbol{\theta}}
\newcommand{\bfmu}{\boldsymbol{\mu}}
\newcommand{\bflambda}{\boldsymbol{\lambda}}
\newcommand{\bfOmega}{\boldsymbol{\Omega}}
\newcommand{\bfSigma}{\boldsymbol{\Sigma}}
\newcommand{\bfPhi}{\boldsymbol{\Phi}}
\newcommand{\bfPsi}{\boldsymbol{\Psi}}
\newcommand{\bfeps}{\boldsymbol{\varepsilon}}
\newcommand{\bfeta}{\boldsymbol{\eta}}
\newcommand{\bfzero}{\mathbf{0}}
\newcommand{\bfpi}{\boldsymbol{\pi}}
\newcommand{\bfrho}{\boldsymbol{\rho}}
\newcommand{\bfLambda}{\boldsymbol{\Lambda}}
\newcommand{\bfv}{\mathbf{v}}

% Epsilon shorthand
\newcommand{\eps}{\varepsilon}
